\section{Learning Objectives}

After completing this chapter, readers will be able to:
\begin{itemize}
    \item Understand the cost structure of ML systems across training and serving
    \item Implement comprehensive training cost optimization strategies
    \item Analyze and reduce serving infrastructure costs
    \item Design intelligent resource allocation and scaling policies
    \item Deploy cloud cost monitoring systems with automated alerts
    \item Leverage spot instances and preemptible VMs effectively
    \item Calculate and optimize ROI for ML investments
    \item Apply cost optimization best practices across the ML lifecycle
\end{itemize}

\section{Introduction}

ML systems often represent significant infrastructure investments, with costs scaling rapidly as model complexity and usage increase. A production ML system can cost anywhere from thousands to millions of dollars annually, making cost optimization a critical engineering concern.

\textbf{Primary Cost Centers in ML Systems:}

\begin{itemize}
    \item \textbf{Training Costs:} Compute for model training, hyperparameter optimization, experimentation
    \item \textbf{Serving Costs:} Inference infrastructure, API endpoints, model hosting
    \item \textbf{Storage Costs:} Training data, model artifacts, feature stores, logs
    \item \textbf{Data Processing:} ETL pipelines, feature engineering, data validation
    \item \textbf{Personnel:} ML engineers, data scientists, MLOps specialists
    \item \textbf{Monitoring and Observability:} Metrics, logs, traces, dashboards
    \item \textbf{Development and Experimentation:} Notebooks, development environments, testing
\end{itemize}

\textbf{Cost Optimization Principles:}

\begin{enumerate}
    \item \textbf{Visibility:} You can't optimize what you can't measure
    \item \textbf{Right-sizing:} Match resources to actual requirements
    \item \textbf{Automation:} Scale resources dynamically based on demand
    \item \textbf{Efficiency:} Optimize algorithms and infrastructure utilization
    \item \textbf{Trade-offs:} Balance cost, performance, and latency requirements
\end{enumerate}

\section{ML Cost Structure and Analysis}

Understanding the cost breakdown of ML systems is the first step toward optimization.

\subsection{Cost Breakdown Model}

\begin{lstlisting}[language=Python, caption=ML System Cost Analyzer]
#!/usr/bin/env python3
"""
ML System Cost Analysis Framework
Comprehensive cost tracking and breakdown for ML infrastructure
"""

from dataclasses import dataclass, field
from typing import Dict, List, Optional
from datetime import datetime, timedelta
from enum import Enum
import pandas as pd
import numpy as np

class CostCategory(Enum):
    """Cost categories for ML systems"""
    TRAINING = "training"
    SERVING = "serving"
    STORAGE = "storage"
    DATA_PROCESSING = "data_processing"
    MONITORING = "monitoring"
    DEVELOPMENT = "development"
    NETWORKING = "networking"

@dataclass
class ResourceCost:
    """Individual resource cost record"""
    resource_id: str
    resource_type: str  # compute, storage, network
    category: CostCategory
    cost: float
    unit: str  # per hour, per GB, per request
    quantity: float
    start_time: datetime
    end_time: datetime
    metadata: Dict = field(default_factory=dict)

    @property
    def total_cost(self) -> float:
        """Calculate total cost"""
        return self.cost * self.quantity

    @property
    def duration_hours(self) -> float:
        """Calculate duration in hours"""
        return (self.end_time - self.start_time).total_seconds() / 3600

@dataclass
class CostReport:
    """Cost analysis report"""
    total_cost: float
    breakdown_by_category: Dict[CostCategory, float]
    breakdown_by_resource_type: Dict[str, float]
    top_cost_drivers: List[tuple]
    optimization_opportunities: List[str]
    period_start: datetime
    period_end: datetime

class MLCostAnalyzer:
    """Analyze and track ML system costs"""

    def __init__(self):
        self.cost_records: List[ResourceCost] = []
        self.hourly_rates = {
            # AWS pricing (example rates)
            "ml.p3.2xlarge": 3.06,  # V100 GPU
            "ml.p4d.24xlarge": 32.77,  # A100 GPU
            "ml.g4dn.xlarge": 0.526,  # T4 GPU
            "ml.m5.xlarge": 0.192,  # General compute
            "ml.c5.2xlarge": 0.34,  # Compute optimized
            # Storage (per GB-month)
            "s3_standard": 0.023,
            "ebs_gp3": 0.08,
            "efs": 0.30,
        }

    def add_cost_record(self, record: ResourceCost):
        """Add a cost record"""
        self.cost_records.append(record)

    def track_training_run(
        self,
        job_id: str,
        instance_type: str,
        num_instances: int,
        start_time: datetime,
        end_time: datetime,
        metadata: Optional[Dict] = None
    ) -> ResourceCost:
        """Track cost of a training run"""

        duration_hours = (end_time - start_time).total_seconds() / 3600
        hourly_cost = self.hourly_rates.get(instance_type, 0)
        total_quantity = duration_hours * num_instances

        record = ResourceCost(
            resource_id=job_id,
            resource_type="compute",
            category=CostCategory.TRAINING,
            cost=hourly_cost,
            unit="per_hour",
            quantity=total_quantity,
            start_time=start_time,
            end_time=end_time,
            metadata=metadata or {}
        )

        self.add_cost_record(record)
        return record

    def track_serving_endpoint(
        self,
        endpoint_name: str,
        instance_type: str,
        num_instances: int,
        requests: int,
        start_time: datetime,
        end_time: datetime
    ) -> ResourceCost:
        """Track cost of serving endpoint"""

        duration_hours = (end_time - start_time).total_seconds() / 3600
        hourly_cost = self.hourly_rates.get(instance_type, 0)
        total_quantity = duration_hours * num_instances

        record = ResourceCost(
            resource_id=endpoint_name,
            resource_type="compute",
            category=CostCategory.SERVING,
            cost=hourly_cost,
            unit="per_hour",
            quantity=total_quantity,
            start_time=start_time,
            end_time=end_time,
            metadata={
                "requests": requests,
                "cost_per_request": (hourly_cost * total_quantity) / requests if requests > 0 else 0
            }
        )

        self.add_cost_record(record)
        return record

    def track_storage(
        self,
        storage_id: str,
        storage_type: str,
        size_gb: float,
        duration_days: int,
        category: CostCategory
    ) -> ResourceCost:
        """Track storage costs"""

        monthly_cost_per_gb = self.hourly_rates.get(storage_type, 0)
        quantity = size_gb * (duration_days / 30.0)  # GB-months

        start_time = datetime.now()
        end_time = start_time + timedelta(days=duration_days)

        record = ResourceCost(
            resource_id=storage_id,
            resource_type="storage",
            category=category,
            cost=monthly_cost_per_gb,
            unit="per_gb_month",
            quantity=quantity,
            start_time=start_time,
            end_time=end_time,
            metadata={"size_gb": size_gb}
        )

        self.add_cost_record(record)
        return record

    def generate_cost_report(
        self,
        start_date: datetime,
        end_date: datetime
    ) -> CostReport:
        """Generate comprehensive cost report"""

        # Filter records by date range
        filtered_records = [
            r for r in self.cost_records
            if r.start_time >= start_date and r.end_time <= end_date
        ]

        # Calculate total cost
        total_cost = sum(r.total_cost for r in filtered_records)

        # Breakdown by category
        breakdown_by_category = {}
        for category in CostCategory:
            category_records = [r for r in filtered_records if r.category == category]
            breakdown_by_category[category] = sum(r.total_cost for r in category_records)

        # Breakdown by resource type
        breakdown_by_resource_type = {}
        for record in filtered_records:
            if record.resource_type not in breakdown_by_resource_type:
                breakdown_by_resource_type[record.resource_type] = 0
            breakdown_by_resource_type[record.resource_type] += record.total_cost

        # Identify top cost drivers
        cost_by_resource = {}
        for record in filtered_records:
            if record.resource_id not in cost_by_resource:
                cost_by_resource[record.resource_id] = 0
            cost_by_resource[record.resource_id] += record.total_cost

        top_cost_drivers = sorted(
            cost_by_resource.items(),
            key=lambda x: x[1],
            reverse=True
        )[:10]

        # Identify optimization opportunities
        optimization_opportunities = self._identify_optimization_opportunities(
            filtered_records
        )

        return CostReport(
            total_cost=total_cost,
            breakdown_by_category=breakdown_by_category,
            breakdown_by_resource_type=breakdown_by_resource_type,
            top_cost_drivers=top_cost_drivers,
            optimization_opportunities=optimization_opportunities,
            period_start=start_date,
            period_end=end_date
        )

    def _identify_optimization_opportunities(
        self,
        records: List[ResourceCost]
    ) -> List[str]:
        """Identify cost optimization opportunities"""

        opportunities = []

        # Check for long-running training jobs
        training_records = [r for r in records if r.category == CostCategory.TRAINING]
        for record in training_records:
            if record.duration_hours > 24:
                opportunities.append(
                    f"Long training job {record.resource_id} ran for "
                    f"{record.duration_hours:.1f} hours. Consider checkpointing "
                    "and spot instances."
                )

        # Check for underutilized serving endpoints
        serving_records = [r for r in records if r.category == CostCategory.SERVING]
        for record in serving_records:
            requests = record.metadata.get("requests", 0)
            if requests > 0:
                cost_per_1k_requests = (record.total_cost / requests) * 1000
                if cost_per_1k_requests > 1.0:
                    opportunities.append(
                        f"Serving endpoint {record.resource_id} has high cost per request "
                        f"(${cost_per_1k_requests:.2f} per 1K). Consider right-sizing or serverless."
                    )

        # Check for expensive storage
        storage_records = [r for r in records if r.resource_type == "storage"]
        for record in storage_records:
            size_gb = record.metadata.get("size_gb", 0)
            if size_gb > 1000 and record.cost > 0.1:  # >1TB on expensive storage
                opportunities.append(
                    f"Large dataset {record.resource_id} ({size_gb:.0f} GB) on expensive "
                    "storage. Consider tiering to cheaper storage class."
                )

        return opportunities

    def export_to_dataframe(self) -> pd.DataFrame:
        """Export cost records to pandas DataFrame"""

        data = []
        for record in self.cost_records:
            data.append({
                "resource_id": record.resource_id,
                "resource_type": record.resource_type,
                "category": record.category.value,
                "cost_per_unit": record.cost,
                "quantity": record.quantity,
                "total_cost": record.total_cost,
                "start_time": record.start_time,
                "end_time": record.end_time,
                "duration_hours": record.duration_hours
            })

        return pd.DataFrame(data)

    def calculate_monthly_projection(self, days_elapsed: int) -> float:
        """Project monthly cost based on current spending"""

        if days_elapsed == 0:
            return 0

        current_cost = sum(r.total_cost for r in self.cost_records)
        daily_average = current_cost / days_elapsed
        monthly_projection = daily_average * 30

        return monthly_projection

# Example usage
if __name__ == "__main__":
    analyzer = MLCostAnalyzer()

    # Track training run
    start = datetime.now() - timedelta(hours=5)
    end = datetime.now()

    analyzer.track_training_run(
        job_id="bert-training-001",
        instance_type="ml.p3.2xlarge",
        num_instances=4,
        start_time=start,
        end_time=end,
        metadata={"model": "BERT", "dataset": "wikipedia"}
    )

    # Track serving endpoint
    analyzer.track_serving_endpoint(
        endpoint_name="recommendation-api",
        instance_type="ml.m5.xlarge",
        num_instances=3,
        requests=1500000,
        start_time=start,
        end_time=end
    )

    # Generate report
    report = analyzer.generate_cost_report(
        start_date=start - timedelta(days=30),
        end_date=end
    )

    print(f"Total Cost: ${report.total_cost:.2f}")
    print("\nCost Breakdown:")
    for category, cost in report.breakdown_by_category.items():
        if cost > 0:
            percentage = (cost / report.total_cost) * 100
            print(f"  {category.value}: ${cost:.2f} ({percentage:.1f}%)")

    print("\nOptimization Opportunities:")
    for opportunity in report.optimization_opportunities:
        print(f"  - {opportunity}")
\end{lstlisting}

\section{Training Cost Optimization}

Training costs often dominate ML budgets, especially for large models and extensive hyperparameter searches.

\subsection{Training Infrastructure Optimization}

\begin{bestpractice}
Profile training jobs to identify compute bottlenecks before scaling to expensive hardware. A poorly optimized training loop won't benefit from larger instances.
\end{bestpractice}

\textbf{Key Training Cost Optimization Strategies:}

\begin{enumerate}
    \item \textbf{Right-size Compute Resources:} Match instance types to workload requirements
    \item \textbf{Use Mixed Precision Training:} Reduce memory and increase throughput with FP16/BF16
    \item \textbf{Implement Gradient Accumulation:} Train larger models on smaller instances
    \item \textbf{Optimize Data Loading:} Eliminate I/O bottlenecks that waste GPU cycles
    \item \textbf{Use Spot Instances:} Save 70-90\% on compute for fault-tolerant workloads
    \item \textbf{Implement Early Stopping:} Don't waste resources on unpromising runs
    \item \textbf{Cache Preprocessed Data:} Avoid redundant data transformations
\end{enumerate}

\subsection{Training Cost Optimizer Implementation}

\begin{lstlisting}[language=Python, caption=Training Cost Optimization Framework]
#!/usr/bin/env python3
"""
Training Cost Optimization
Strategies and tools for reducing ML training costs
"""

import torch
import torch.nn as nn
import time
from typing import Optional, Dict, List, Tuple
from dataclasses import dataclass
import psutil
import GPUtil

@dataclass
class TrainingMetrics:
    """Training performance metrics"""
    samples_per_second: float
    gpu_utilization: float
    memory_utilization: float
    cost_per_hour: float
    estimated_total_cost: float
    time_to_completion_hours: float

    @property
    def cost_per_1k_samples(self) -> float:
        """Calculate cost per 1000 training samples"""
        if self.samples_per_second == 0:
            return float('inf')
        samples_per_hour = self.samples_per_second * 3600
        return (self.cost_per_hour / samples_per_hour) * 1000

class TrainingCostOptimizer:
    """Optimize training costs through various strategies"""

    def __init__(
        self,
        model: nn.Module,
        cost_per_hour: float,
        target_samples: int
    ):
        self.model = model
        self.cost_per_hour = cost_per_hour
        self.target_samples = target_samples
        self.training_start_time = None
        self.samples_processed = 0

    def profile_training_step(
        self,
        batch_size: int,
        num_steps: int = 100
    ) -> TrainingMetrics:
        """Profile training performance and cost"""

        device = next(self.model.parameters()).device
        dummy_input = torch.randn(batch_size, 3, 224, 224).to(device)
        dummy_target = torch.randint(0, 1000, (batch_size,)).to(device)

        # Warmup
        for _ in range(10):
            output = self.model(dummy_input)
            loss = nn.functional.cross_entropy(output, dummy_target)
            loss.backward()

        # Profile
        start_time = time.time()

        for _ in range(num_steps):
            output = self.model(dummy_input)
            loss = nn.functional.cross_entropy(output, dummy_target)
            loss.backward()

        torch.cuda.synchronize() if torch.cuda.is_available() else None
        elapsed_time = time.time() - start_time

        # Calculate metrics
        samples_per_second = (num_steps * batch_size) / elapsed_time

        # Get GPU utilization
        gpu_utilization = 0
        memory_utilization = 0
        if torch.cuda.is_available():
            gpus = GPUtil.getGPUs()
            if gpus:
                gpu_utilization = gpus[0].load * 100
                memory_utilization = gpus[0].memoryUtil * 100

        # Cost calculations
        remaining_samples = self.target_samples - self.samples_processed
        time_to_completion_hours = remaining_samples / (samples_per_second * 3600)
        estimated_total_cost = time_to_completion_hours * self.cost_per_hour

        return TrainingMetrics(
            samples_per_second=samples_per_second,
            gpu_utilization=gpu_utilization,
            memory_utilization=memory_utilization,
            cost_per_hour=self.cost_per_hour,
            estimated_total_cost=estimated_total_cost,
            time_to_completion_hours=time_to_completion_hours
        )

    def find_optimal_batch_size(
        self,
        min_batch_size: int = 8,
        max_batch_size: int = 512
    ) -> Tuple[int, TrainingMetrics]:
        """Find batch size that maximizes throughput per dollar"""

        best_batch_size = min_batch_size
        best_cost_efficiency = float('inf')
        best_metrics = None

        batch_size = min_batch_size
        while batch_size <= max_batch_size:
            try:
                metrics = self.profile_training_step(batch_size)

                # Cost efficiency: dollars per 1000 samples
                cost_efficiency = metrics.cost_per_1k_samples

                if cost_efficiency < best_cost_efficiency:
                    best_cost_efficiency = cost_efficiency
                    best_batch_size = batch_size
                    best_metrics = metrics

                batch_size *= 2

            except RuntimeError as e:
                # OOM error, stop increasing batch size
                if "out of memory" in str(e):
                    break
                raise

        return best_batch_size, best_metrics

    def enable_mixed_precision(self) -> Dict[str, float]:
        """Enable mixed precision training for cost savings"""

        # Profile FP32
        metrics_fp32 = self.profile_training_step(batch_size=32)

        # Enable mixed precision
        scaler = torch.cuda.amp.GradScaler()

        # Profile with mixed precision
        device = next(self.model.parameters()).device
        dummy_input = torch.randn(32, 3, 224, 224).to(device)
        dummy_target = torch.randint(0, 1000, (32,)).to(device)

        start_time = time.time()
        for _ in range(100):
            with torch.cuda.amp.autocast():
                output = self.model(dummy_input)
                loss = nn.functional.cross_entropy(output, dummy_target)

            scaler.scale(loss).backward()
            scaler.step(torch.optim.SGD(self.model.parameters(), lr=0.01))
            scaler.update()

        torch.cuda.synchronize() if torch.cuda.is_available() else None
        elapsed_time = time.time() - start_time

        samples_per_second_fp16 = (100 * 32) / elapsed_time

        # Calculate speedup and cost savings
        speedup = samples_per_second_fp16 / metrics_fp32.samples_per_second
        cost_reduction = (1 - (1 / speedup)) * 100

        return {
            "speedup": speedup,
            "cost_reduction_percent": cost_reduction,
            "fp32_samples_per_sec": metrics_fp32.samples_per_second,
            "fp16_samples_per_sec": samples_per_second_fp16
        }

    def implement_gradient_accumulation(
        self,
        target_batch_size: int,
        available_memory_batch_size: int
    ) -> Dict[str, any]:
        """Use gradient accumulation to simulate larger batches"""

        accumulation_steps = target_batch_size // available_memory_batch_size

        # This allows training with larger effective batch sizes
        # on smaller, cheaper instances

        # Cost comparison
        large_instance_cost = self.cost_per_hour * 2  # Assume 2x cost
        small_instance_time_multiplier = 1.1  # Slight overhead from accumulation

        large_instance_total = large_instance_cost
        small_instance_total = self.cost_per_hour * small_instance_time_multiplier

        cost_savings = ((large_instance_total - small_instance_total) /
                       large_instance_total) * 100

        return {
            "accumulation_steps": accumulation_steps,
            "effective_batch_size": target_batch_size,
            "physical_batch_size": available_memory_batch_size,
            "estimated_cost_savings_percent": cost_savings
        }

    def calculate_early_stopping_savings(
        self,
        total_epochs: int,
        early_stop_epoch: int,
        num_runs: int
    ) -> Dict[str, float]:
        """Calculate cost savings from early stopping"""

        # Without early stopping
        cost_per_epoch = self.cost_per_hour * (self.target_samples /
                        (self.profile_training_step(32).samples_per_second * 3600))
        total_cost_without = cost_per_epoch * total_epochs * num_runs

        # With early stopping (average case)
        average_stop_epoch = (early_stop_epoch + total_epochs) / 2
        total_cost_with = cost_per_epoch * average_stop_epoch * num_runs

        savings = total_cost_without - total_cost_with
        savings_percent = (savings / total_cost_without) * 100

        return {
            "cost_without_early_stopping": total_cost_without,
            "cost_with_early_stopping": total_cost_with,
            "savings_dollars": savings,
            "savings_percent": savings_percent
        }

class HyperparameterSearchCostOptimizer:
    """Optimize costs for hyperparameter search"""

    def __init__(self, cost_per_hour: float):
        self.cost_per_hour = cost_per_hour

    def compare_search_strategies(
        self,
        num_hyperparameters: int,
        values_per_param: int,
        max_trials: int,
        avg_trial_hours: float
    ) -> Dict[str, Dict[str, float]]:
        """Compare cost of different search strategies"""

        # Grid Search
        grid_search_trials = values_per_param ** num_hyperparameters
        grid_search_cost = min(grid_search_trials, max_trials) * avg_trial_hours * self.cost_per_hour

        # Random Search
        random_search_trials = max_trials
        random_search_cost = random_search_trials * avg_trial_hours * self.cost_per_hour

        # Bayesian Optimization (typically needs fewer trials)
        bayesian_trials = max_trials // 2
        bayesian_cost = bayesian_trials * avg_trial_hours * self.cost_per_hour

        # Hyperband (early stopping poor configurations)
        hyperband_trials = max_trials
        hyperband_avg_hours = avg_trial_hours * 0.3  # Many trials stopped early
        hyperband_cost = hyperband_trials * hyperband_avg_hours * self.cost_per_hour

        return {
            "grid_search": {
                "trials": grid_search_trials,
                "cost": grid_search_cost,
                "hours": grid_search_trials * avg_trial_hours
            },
            "random_search": {
                "trials": random_search_trials,
                "cost": random_search_cost,
                "hours": random_search_trials * avg_trial_hours
            },
            "bayesian_optimization": {
                "trials": bayesian_trials,
                "cost": bayesian_cost,
                "hours": bayesian_trials * avg_trial_hours
            },
            "hyperband": {
                "trials": hyperband_trials,
                "cost": hyperband_cost,
                "hours": hyperband_trials * hyperband_avg_hours
            }
        }

    def calculate_sequential_halving_savings(
        self,
        initial_configurations: int,
        max_budget_per_config: float,
        halving_rounds: int
    ) -> Dict[str, float]:
        """Calculate savings from successive halving"""

        # Full training cost
        full_cost = initial_configurations * max_budget_per_config * self.cost_per_hour

        # Successive halving cost
        halving_cost = 0
        configs = initial_configurations
        budget = max_budget_per_config / (2 ** (halving_rounds - 1))

        for round_num in range(halving_rounds):
            halving_cost += configs * budget * self.cost_per_hour
            configs = configs // 2
            budget = budget * 2

        savings = full_cost - halving_cost
        savings_percent = (savings / full_cost) * 100

        return {
            "full_training_cost": full_cost,
            "successive_halving_cost": halving_cost,
            "savings_dollars": savings,
            "savings_percent": savings_percent
        }

# Example usage
if __name__ == "__main__":
    # Create dummy model
    model = torch.nn.Sequential(
        torch.nn.Linear(100, 256),
        torch.nn.ReLU(),
        torch.nn.Linear(256, 10)
    )

    if torch.cuda.is_available():
        model = model.cuda()

    # Initialize optimizer
    optimizer = TrainingCostOptimizer(
        model=model,
        cost_per_hour=3.06,  # ml.p3.2xlarge
        target_samples=1000000
    )

    # Find optimal batch size
    optimal_batch, metrics = optimizer.find_optimal_batch_size()
    print(f"Optimal batch size: {optimal_batch}")
    print(f"Cost per 1K samples: ${metrics.cost_per_1k_samples:.4f}")
    print(f"Estimated total cost: ${metrics.estimated_total_cost:.2f}")

    # Analyze mixed precision savings
    mp_savings = optimizer.enable_mixed_precision()
    print(f"\nMixed Precision Speedup: {mp_savings['speedup']:.2f}x")
    print(f"Cost Reduction: {mp_savings['cost_reduction_percent']:.1f}%")
\end{lstlisting}

\subsection{Spot Instance Strategy for Training}

Spot instances can reduce training costs by 70-90\%, but require fault-tolerance mechanisms.

\begin{lstlisting}[language=Python, caption=Spot Instance Training Manager]
#!/usr/bin/env python3
"""
Spot Instance Training Management
Handle interruptions and optimize costs with spot instances
"""

import boto3
import time
import json
from datetime import datetime
from typing import Optional, Dict, List
from dataclasses import dataclass
import threading

@dataclass
class SpotInstanceConfig:
    """Spot instance configuration"""
    instance_type: str
    max_price: float  # Maximum price per hour
    availability_zones: List[str]
    interruption_behavior: str = "terminate"  # or "stop", "hibernate"

@dataclass
class CheckpointConfig:
    """Checkpointing configuration"""
    checkpoint_interval_minutes: int
    s3_bucket: str
    checkpoint_path: str
    max_checkpoints_to_keep: int = 3

class SpotInstanceTrainingManager:
    """Manage training on spot instances with fault tolerance"""

    def __init__(
        self,
        spot_config: SpotInstanceConfig,
        checkpoint_config: CheckpointConfig,
        aws_region: str = "us-east-1"
    ):
        self.spot_config = spot_config
        self.checkpoint_config = checkpoint_config
        self.ec2_client = boto3.client('ec2', region_name=aws_region)
        self.s3_client = boto3.client('s3', region_name=aws_region)
        self.instance_id: Optional[str] = None
        self.interruption_monitor_thread: Optional[threading.Thread] = None
        self.interrupted = False

    def request_spot_instance(self) -> str:
        """Request a spot instance"""

        # Get current spot price
        current_price = self._get_current_spot_price()

        if current_price > self.spot_config.max_price:
            raise ValueError(
                f"Current spot price (${current_price:.3f}) exceeds "
                f"max price (${self.spot_config.max_price:.3f})"
            )

        # Create launch specification
        launch_specification = {
            'ImageId': 'ami-0c55b159cbfafe1f0',  # Deep Learning AMI
            'InstanceType': self.spot_config.instance_type,
            'KeyName': 'your-key-pair',
            'SecurityGroups': ['ml-training-sg'],
            'UserData': self._generate_user_data(),
            'IamInstanceProfile': {
                'Name': 'ml-training-role'
            },
            'BlockDeviceMappings': [
                {
                    'DeviceName': '/dev/sda1',
                    'Ebs': {
                        'VolumeSize': 100,
                        'VolumeType': 'gp3'
                    }
                }
            ]
        }

        # Request spot instance
        response = self.ec2_client.request_spot_instances(
            SpotPrice=str(self.spot_config.max_price),
            InstanceCount=1,
            Type='one-time',
            LaunchSpecification=launch_specification,
            InstanceInterruptionBehavior=self.spot_config.interruption_behavior
        )

        request_id = response['SpotInstanceRequests'][0]['SpotInstanceRequestId']

        # Wait for instance to be running
        self.instance_id = self._wait_for_spot_instance(request_id)

        # Start monitoring for interruptions
        self._start_interruption_monitoring()

        return self.instance_id

    def _get_current_spot_price(self) -> float:
        """Get current spot price for instance type"""

        response = self.ec2_client.describe_spot_price_history(
            InstanceTypes=[self.spot_config.instance_type],
            MaxResults=1,
            ProductDescriptions=['Linux/UNIX'],
            AvailabilityZone=self.spot_config.availability_zones[0]
        )

        if response['SpotPriceHistory']:
            return float(response['SpotPriceHistory'][0]['SpotPrice'])
        return float('inf')

    def _wait_for_spot_instance(self, request_id: str, timeout: int = 300) -> str:
        """Wait for spot instance to be fulfilled"""

        start_time = time.time()

        while time.time() - start_time < timeout:
            response = self.ec2_client.describe_spot_instance_requests(
                SpotInstanceRequestIds=[request_id]
            )

            request = response['SpotInstanceRequests'][0]
            state = request['State']

            if state == 'active':
                return request['InstanceId']
            elif state in ['closed', 'cancelled', 'failed']:
                raise RuntimeError(f"Spot request failed: {request.get('Fault', {})}")

            time.sleep(10)

        raise TimeoutError("Spot instance request timed out")

    def _start_interruption_monitoring(self):
        """Start monitoring for spot interruption notices"""

        def monitor():
            # Check EC2 instance metadata for interruption notice
            # This would run on the instance itself
            while not self.interrupted:
                try:
                    # Check for interruption notice
                    # (simplified - actual implementation would check instance metadata)
                    time.sleep(5)
                except Exception:
                    pass

        self.interruption_monitor_thread = threading.Thread(target=monitor)
        self.interruption_monitor_thread.daemon = True
        self.interruption_monitor_thread.start()

    def save_checkpoint(
        self,
        model_state: Dict,
        optimizer_state: Dict,
        epoch: int,
        global_step: int,
        metadata: Optional[Dict] = None
    ):
        """Save training checkpoint to S3"""

        checkpoint = {
            'model_state_dict': model_state,
            'optimizer_state_dict': optimizer_state,
            'epoch': epoch,
            'global_step': global_step,
            'timestamp': datetime.now().isoformat(),
            'metadata': metadata or {}
        }

        # Save locally first
        local_path = f'/tmp/checkpoint_{global_step}.pt'
        import torch
        torch.save(checkpoint, local_path)

        # Upload to S3
        s3_key = f"{self.checkpoint_config.checkpoint_path}/checkpoint_{global_step}.pt"

        self.s3_client.upload_file(
            local_path,
            self.checkpoint_config.s3_bucket,
            s3_key
        )

        print(f"Checkpoint saved to s3://{self.checkpoint_config.s3_bucket}/{s3_key}")

        # Clean up old checkpoints
        self._cleanup_old_checkpoints()

    def load_latest_checkpoint(self) -> Optional[Dict]:
        """Load the most recent checkpoint from S3"""

        # List checkpoints
        response = self.s3_client.list_objects_v2(
            Bucket=self.checkpoint_config.s3_bucket,
            Prefix=self.checkpoint_config.checkpoint_path
        )

        if 'Contents' not in response:
            return None

        # Sort by last modified
        checkpoints = sorted(
            response['Contents'],
            key=lambda x: x['LastModified'],
            reverse=True
        )

        if not checkpoints:
            return None

        # Download latest checkpoint
        latest_key = checkpoints[0]['Key']
        local_path = '/tmp/latest_checkpoint.pt'

        self.s3_client.download_file(
            self.checkpoint_config.s3_bucket,
            latest_key,
            local_path
        )

        import torch
        checkpoint = torch.load(local_path)

        print(f"Loaded checkpoint from epoch {checkpoint['epoch']}, "
              f"step {checkpoint['global_step']}")

        return checkpoint

    def _cleanup_old_checkpoints(self):
        """Remove old checkpoints to save storage costs"""

        response = self.s3_client.list_objects_v2(
            Bucket=self.checkpoint_config.s3_bucket,
            Prefix=self.checkpoint_config.checkpoint_path
        )

        if 'Contents' not in response:
            return

        checkpoints = sorted(
            response['Contents'],
            key=lambda x: x['LastModified'],
            reverse=True
        )

        # Delete old checkpoints
        checkpoints_to_delete = checkpoints[self.checkpoint_config.max_checkpoints_to_keep:]

        for checkpoint in checkpoints_to_delete:
            self.s3_client.delete_object(
                Bucket=self.checkpoint_config.s3_bucket,
                Key=checkpoint['Key']
            )

    def _generate_user_data(self) -> str:
        """Generate user data script for instance startup"""

        script = f"""#!/bin/bash
# Download training code and resume from checkpoint
aws s3 cp s3://{self.checkpoint_config.s3_bucket}/training_code/ /home/ubuntu/training/ --recursive

# Start training with automatic checkpoint resumption
cd /home/ubuntu/training
python train.py --resume-from-checkpoint
"""
        return script

    def calculate_cost_savings(
        self,
        training_hours: float,
        on_demand_price: float,
        interruption_overhead_percent: float = 5.0
    ) -> Dict[str, float]:
        """Calculate cost savings from using spot instances"""

        # On-demand cost
        on_demand_cost = training_hours * on_demand_price

        # Spot cost (including interruption overhead)
        effective_training_hours = training_hours * (1 + interruption_overhead_percent / 100)
        spot_price = self._get_current_spot_price()
        spot_cost = effective_training_hours * spot_price

        # Savings
        savings = on_demand_cost - spot_cost
        savings_percent = (savings / on_demand_cost) * 100

        return {
            "on_demand_cost": on_demand_cost,
            "spot_cost": spot_cost,
            "savings_dollars": savings,
            "savings_percent": savings_percent,
            "spot_price": spot_price,
            "on_demand_price": on_demand_price
        }

# Example usage
if __name__ == "__main__":
    spot_config = SpotInstanceConfig(
        instance_type="ml.p3.2xlarge",
        max_price=1.50,  # 50% of on-demand price
        availability_zones=["us-east-1a", "us-east-1b"]
    )

    checkpoint_config = CheckpointConfig(
        checkpoint_interval_minutes=15,
        s3_bucket="my-ml-checkpoints",
        checkpoint_path="training/resnet50"
    )

    manager = SpotInstanceTrainingManager(spot_config, checkpoint_config)

    # Calculate potential savings
    savings = manager.calculate_cost_savings(
        training_hours=48,
        on_demand_price=3.06
    )

    print(f"Training on spot instances:")
    print(f"  On-demand cost: ${savings['on_demand_cost']:.2f}")
    print(f"  Spot cost: ${savings['spot_cost']:.2f}")
    print(f"  Savings: ${savings['savings_dollars']:.2f} ({savings['savings_percent']:.1f}%)")
\end{lstlisting}

\section{Serving Cost Analysis and Reduction}

Serving costs scale with traffic and can exceed training costs for high-traffic applications.

\subsection{Serving Cost Optimization Strategies}

\textbf{Key Serving Cost Drivers:}

\begin{itemize}
    \item \textbf{Compute Resources:} Instance types and count for model hosting
    \item \textbf{Request Volume:} Number of inference requests
    \item \textbf{Model Complexity:} Inference latency and resource requirements
    \item \textbf{Always-On Infrastructure:} Fixed costs vs. usage-based costs
    \item \textbf{Data Transfer:} Network egress charges
\end{itemize}

\begin{lstlisting}[language=Python, caption=Serving Cost Analyzer and Optimizer]
#!/usr/bin/env python3
"""
Serving Cost Optimization
Analyze and optimize model serving infrastructure costs
"""

from dataclasses import dataclass
from typing import Dict, List, Optional
import numpy as np

@dataclass
class ServingConfig:
    """Model serving configuration"""
    instance_type: str
    num_instances: int
    requests_per_second: float
    avg_latency_ms: float
    cost_per_hour: float

    @property
    def total_hourly_cost(self) -> float:
        return self.cost_per_hour * self.num_instances

    @property
    def cost_per_1m_requests(self) -> float:
        requests_per_hour = self.requests_per_second * 3600
        if requests_per_hour == 0:
            return float('inf')
        return (self.total_hourly_cost / requests_per_hour) * 1_000_000

@dataclass
class ServerlessConfig:
    """Serverless serving configuration"""
    cost_per_request: float
    cost_per_gb_second: float
    avg_memory_gb: float
    avg_duration_seconds: float

    def cost_per_1m_requests(self) -> float:
        compute_cost = self.cost_per_gb_second * self.avg_memory_gb * self.avg_duration_seconds
        total_cost_per_request = self.cost_per_request + compute_cost
        return total_cost_per_request * 1_000_000

class ServingCostOptimizer:
    """Optimize model serving costs"""

    def __init__(self):
        self.pricing = {
            # AWS pricing (example)
            "ml.m5.xlarge": 0.269,
            "ml.m5.2xlarge": 0.538,
            "ml.c5.xlarge": 0.238,
            "ml.c5.2xlarge": 0.476,
            "ml.g4dn.xlarge": 0.736,  # With GPU
            "ml.inf1.xlarge": 0.368,  # AWS Inferentia
        }

    def compare_deployment_options(
        self,
        daily_requests: int,
        avg_latency_ms: float,
        p99_latency_ms: float,
        peak_rps: float,
        avg_rps: float
    ) -> Dict[str, Dict]:
        """Compare different deployment strategies"""

        results = {}

        # Option 1: Always-on instances sized for peak
        peak_instances = self._calculate_required_instances(
            peak_rps,
            instance_type="ml.m5.xlarge",
            requests_per_instance=100
        )

        results["always_on_peak"] = {
            "instance_type": "ml.m5.xlarge",
            "num_instances": peak_instances,
            "monthly_cost": peak_instances * self.pricing["ml.m5.xlarge"] * 24 * 30,
            "cost_per_1m_requests": self._calculate_cost_per_1m_requests(
                peak_instances,
                self.pricing["ml.m5.xlarge"],
                daily_requests
            ),
            "description": "Always-on instances sized for peak traffic"
        }

        # Option 2: Always-on for baseline + auto-scaling
        baseline_instances = self._calculate_required_instances(
            avg_rps,
            instance_type="ml.m5.xlarge",
            requests_per_instance=100
        )

        # Assume auto-scaling covers 20% of time at peak
        auto_scale_instance_hours = (peak_instances - baseline_instances) * 24 * 30 * 0.2
        baseline_instance_hours = baseline_instances * 24 * 30

        results["autoscaling"] = {
            "baseline_instances": baseline_instances,
            "peak_instances": peak_instances,
            "monthly_cost": (
                (baseline_instance_hours + auto_scale_instance_hours) *
                self.pricing["ml.m5.xlarge"]
            ),
            "cost_per_1m_requests": self._calculate_cost_per_1m_requests(
                baseline_instances,
                self.pricing["ml.m5.xlarge"],
                daily_requests,
                autoscaling=True
            ),
            "description": "Baseline instances + auto-scaling"
        }

        # Option 3: Serverless (AWS Lambda / SageMaker Serverless)
        # Lambda pricing: $0.20 per 1M requests + compute time
        avg_duration_seconds = avg_latency_ms / 1000
        memory_gb = 2.0  # Example memory allocation

        compute_cost_per_request = 0.0000166667 * memory_gb * avg_duration_seconds
        request_cost = 0.20 / 1_000_000
        total_cost_per_request = compute_cost_per_request + request_cost

        results["serverless"] = {
            "monthly_cost": total_cost_per_request * daily_requests * 30,
            "cost_per_1m_requests": total_cost_per_request * 1_000_000,
            "memory_gb": memory_gb,
            "avg_duration_seconds": avg_duration_seconds,
            "description": "Serverless compute (pay per request)"
        }

        # Option 4: Optimized instances (Inferentia/GPU)
        opt_instances = self._calculate_required_instances(
            peak_rps,
            instance_type="ml.inf1.xlarge",
            requests_per_instance=200  # Better throughput with optimization
        )

        results["optimized_hardware"] = {
            "instance_type": "ml.inf1.xlarge",
            "num_instances": opt_instances,
            "monthly_cost": opt_instances * self.pricing["ml.inf1.xlarge"] * 24 * 30,
            "cost_per_1m_requests": self._calculate_cost_per_1m_requests(
                opt_instances,
                self.pricing["ml.inf1.xlarge"],
                daily_requests
            ),
            "description": "Optimized hardware (Inferentia)"
        }

        return results

    def _calculate_required_instances(
        self,
        target_rps: float,
        instance_type: str,
        requests_per_instance: float
    ) -> int:
        """Calculate number of instances needed"""
        # Add 20% headroom for spikes
        required = int(np.ceil(target_rps / requests_per_instance * 1.2))
        return max(required, 1)

    def _calculate_cost_per_1m_requests(
        self,
        num_instances: int,
        cost_per_hour: float,
        daily_requests: int,
        autoscaling: bool = False
    ) -> float:
        """Calculate cost per million requests"""
        monthly_requests = daily_requests * 30
        monthly_instance_hours = num_instances * 24 * 30

        if autoscaling:
            monthly_instance_hours *= 0.7  # Assume 30% reduction with autoscaling

        monthly_cost = monthly_instance_hours * cost_per_hour

        if monthly_requests == 0:
            return float('inf')

        return (monthly_cost / monthly_requests) * 1_000_000

    def optimize_batch_size(
        self,
        model_latency_by_batch: Dict[int, float],
        cost_per_hour: float,
        target_p99_latency_ms: float
    ) -> Dict:
        """Find optimal batch size for serving"""

        best_batch_size = 1
        best_cost_per_1k_requests = float('inf')

        results = {}

        for batch_size, latency_ms in model_latency_by_batch.items():
            if latency_ms > target_p99_latency_ms:
                continue  # Violates latency SLA

            # Throughput (requests per hour)
            requests_per_second = batch_size / (latency_ms / 1000)
            requests_per_hour = requests_per_second * 3600

            # Cost per 1K requests
            cost_per_1k = (cost_per_hour / requests_per_hour) * 1000

            results[batch_size] = {
                "latency_ms": latency_ms,
                "requests_per_second": requests_per_second,
                "cost_per_1k_requests": cost_per_1k
            }

            if cost_per_1k < best_cost_per_1k_requests:
                best_cost_per_1k_requests = cost_per_1k
                best_batch_size = batch_size

        return {
            "optimal_batch_size": best_batch_size,
            "cost_per_1k_requests": best_cost_per_1k_requests,
            "all_results": results
        }

    def calculate_model_compression_savings(
        self,
        original_latency_ms: float,
        compressed_latency_ms: float,
        original_memory_gb: float,
        compressed_memory_gb: float,
        daily_requests: int
    ) -> Dict:
        """Calculate savings from model compression"""

        # Original deployment
        original_instances = max(1, int(np.ceil(original_memory_gb / 8)))  # 8GB per instance
        original_cost_per_hour = self.pricing["ml.m5.xlarge"] * original_instances

        # Compressed deployment
        compressed_instances = max(1, int(np.ceil(compressed_memory_gb / 8)))
        compressed_cost_per_hour = self.pricing["ml.m5.xlarge"] * compressed_instances

        # If latency improved, might handle more load per instance
        throughput_improvement = original_latency_ms / compressed_latency_ms

        # Monthly costs
        original_monthly = original_cost_per_hour * 24 * 30
        compressed_monthly = (compressed_cost_per_hour / throughput_improvement) * 24 * 30

        savings = original_monthly - compressed_monthly
        savings_percent = (savings / original_monthly) * 100

        return {
            "original_monthly_cost": original_monthly,
            "compressed_monthly_cost": compressed_monthly,
            "savings_dollars": savings,
            "savings_percent": savings_percent,
            "latency_improvement_percent": ((original_latency_ms - compressed_latency_ms) /
                                           original_latency_ms) * 100,
            "memory_reduction_percent": ((original_memory_gb - compressed_memory_gb) /
                                        original_memory_gb) * 100
        }

    def analyze_caching_opportunity(
        self,
        total_daily_requests: int,
        cache_hit_rate: float,
        cached_request_cost: float,
        uncached_request_cost: float
    ) -> Dict:
        """Calculate savings from response caching"""

        cache_hits = total_daily_requests * cache_hit_rate
        cache_misses = total_daily_requests * (1 - cache_hit_rate)

        # Without caching
        cost_without_cache = total_daily_requests * uncached_request_cost

        # With caching
        cost_with_cache = (
            cache_hits * cached_request_cost +
            cache_misses * uncached_request_cost
        )

        # Cache infrastructure cost (Redis/Memcached)
        cache_infrastructure_cost_daily = 5.0  # Example: $5/day for cache

        total_cost_with_cache = cost_with_cache + cache_infrastructure_cost_daily

        savings = cost_without_cache - total_cost_with_cache
        savings_percent = (savings / cost_without_cache) * 100

        monthly_savings = savings * 30

        return {
            "cost_without_cache_monthly": cost_without_cache * 30,
            "cost_with_cache_monthly": total_cost_with_cache * 30,
            "savings_monthly": monthly_savings,
            "savings_percent": savings_percent,
            "cache_hit_rate": cache_hit_rate,
            "break_even_hit_rate": self._calculate_break_even_hit_rate(
                uncached_request_cost,
                cached_request_cost,
                cache_infrastructure_cost_daily,
                total_daily_requests
            )
        }

    def _calculate_break_even_hit_rate(
        self,
        uncached_cost: float,
        cached_cost: float,
        infrastructure_cost: float,
        daily_requests: int
    ) -> float:
        """Calculate minimum cache hit rate to break even"""
        # Solve: infrastructure_cost = (hit_rate * daily_requests * (uncached_cost - cached_cost))
        cost_saving_per_hit = uncached_cost - cached_cost
        if cost_saving_per_hit <= 0:
            return 1.0

        break_even_rate = infrastructure_cost / (daily_requests * cost_saving_per_hit)
        return min(break_even_rate, 1.0)

# Example usage
if __name__ == "__main__":
    optimizer = ServingCostOptimizer()

    # Compare deployment options
    comparison = optimizer.compare_deployment_options(
        daily_requests=10_000_000,
        avg_latency_ms=50,
        p99_latency_ms=200,
        peak_rps=500,
        avg_rps=115
    )

    print("Deployment Option Comparison:")
    for option_name, details in comparison.items():
        print(f"\n{option_name}:")
        print(f"  Monthly Cost: ${details['monthly_cost']:.2f}")
        print(f"  Cost per 1M requests: ${details['cost_per_1m_requests']:.2f}")
        print(f"  Description: {details['description']}")

    # Analyze caching
    caching_analysis = optimizer.analyze_caching_opportunity(
        total_daily_requests=10_000_000,
        cache_hit_rate=0.60,
        cached_request_cost=0.001,
        uncached_request_cost=0.01
    )

    print(f"\nCaching Analysis:")
    print(f"  Monthly savings: ${caching_analysis['savings_monthly']:.2f}")
    print(f"  Savings percent: {caching_analysis['savings_percent']:.1f}%")
    print(f"  Break-even hit rate: {caching_analysis['break_even_hit_rate']*100:.1f}%")
\end{lstlisting}

\section{Resource Allocation and Scaling Policies}

Intelligent scaling policies ensure you pay only for resources you need.

\subsection{Autoscaling Implementation}

\begin{lstlisting}[language=Python, caption=ML-Aware Autoscaling Policy]
#!/usr/bin/env python3
"""
ML-Aware Autoscaling
Intelligent scaling policies for ML workloads
"""

from dataclasses import dataclass
from typing import List, Optional, Callable
from datetime import datetime, timedelta
import numpy as np

@dataclass
class ScalingMetrics:
    """Metrics for scaling decisions"""
    timestamp: datetime
    cpu_utilization: float
    memory_utilization: float
    requests_per_second: float
    avg_latency_ms: float
    p99_latency_ms: float
    queue_depth: int

@dataclass
class ScalingPolicy:
    """Autoscaling policy configuration"""
    min_instances: int
    max_instances: int
    target_cpu_utilization: float = 70.0
    target_latency_p99_ms: float = 200.0
    scale_up_cooldown_seconds: int = 300
    scale_down_cooldown_seconds: int = 600
    scale_up_threshold_breaches: int = 2
    scale_down_threshold_breaches: int = 3

class MLAutoscaler:
    """ML-aware autoscaling engine"""

    def __init__(self, policy: ScalingPolicy):
        self.policy = policy
        self.current_instances = policy.min_instances
        self.metrics_history: List[ScalingMetrics] = []
        self.last_scale_up_time: Optional[datetime] = None
        self.last_scale_down_time: Optional[datetime] = None
        self.consecutive_scale_up_signals = 0
        self.consecutive_scale_down_signals = 0

    def add_metrics(self, metrics: ScalingMetrics):
        """Add new metrics for scaling decision"""
        self.metrics_history.append(metrics)

        # Keep only last hour of metrics
        cutoff_time = datetime.now() - timedelta(hours=1)
        self.metrics_history = [
            m for m in self.metrics_history
            if m.timestamp > cutoff_time
        ]

    def should_scale_up(self, metrics: ScalingMetrics) -> bool:
        """Determine if should scale up"""

        if self.current_instances >= self.policy.max_instances:
            return False

        # Check cooldown
        if self.last_scale_up_time:
            elapsed = (datetime.now() - self.last_scale_up_time).total_seconds()
            if elapsed < self.policy.scale_up_cooldown_seconds:
                return False

        # Check thresholds
        should_scale = (
            metrics.cpu_utilization > self.policy.target_cpu_utilization or
            metrics.p99_latency_ms > self.policy.target_latency_p99_ms or
            metrics.queue_depth > 100
        )

        if should_scale:
            self.consecutive_scale_up_signals += 1
            self.consecutive_scale_down_signals = 0
        else:
            self.consecutive_scale_up_signals = 0

        # Require multiple consecutive signals
        return self.consecutive_scale_up_signals >= self.policy.scale_up_threshold_breaches

    def should_scale_down(self, metrics: ScalingMetrics) -> bool:
        """Determine if should scale down"""

        if self.current_instances <= self.policy.min_instances:
            return False

        # Check cooldown
        if self.last_scale_down_time:
            elapsed = (datetime.now() - self.last_scale_down_time).total_seconds()
            if elapsed < self.policy.scale_down_cooldown_seconds:
                return False

        # More conservative for scale down
        should_scale = (
            metrics.cpu_utilization < self.policy.target_cpu_utilization * 0.5 and
            metrics.p99_latency_ms < self.policy.target_latency_p99_ms * 0.7 and
            metrics.queue_depth == 0
        )

        if should_scale:
            self.consecutive_scale_down_signals += 1
            self.consecutive_scale_up_signals = 0
        else:
            self.consecutive_scale_down_signals = 0

        return self.consecutive_scale_down_signals >= self.policy.scale_down_threshold_breaches

    def calculate_target_instances(self, metrics: ScalingMetrics) -> int:
        """Calculate target number of instances"""

        # Calculate based on different metrics
        cpu_based = int(np.ceil(
            self.current_instances *
            (metrics.cpu_utilization / self.policy.target_cpu_utilization)
        ))

        # Estimate instances needed for latency SLA
        # Simple heuristic: if latency is high, we need more capacity
        latency_ratio = metrics.p99_latency_ms / self.policy.target_latency_p99_ms
        latency_based = int(np.ceil(self.current_instances * latency_ratio))

        # Take the maximum
        target = max(cpu_based, latency_based)

        # Constrain to policy limits
        target = max(self.policy.min_instances, min(target, self.policy.max_instances))

        return target

    def execute_scaling_decision(self, metrics: ScalingMetrics) -> Optional[int]:
        """Make and execute scaling decision"""

        self.add_metrics(metrics)

        if self.should_scale_up(metrics):
            target = self.calculate_target_instances(metrics)
            scale_amount = target - self.current_instances

            # Don't scale too aggressively
            scale_amount = min(scale_amount, max(1, self.current_instances // 2))

            new_instances = self.current_instances + scale_amount
            self.current_instances = new_instances
            self.last_scale_up_time = datetime.now()
            self.consecutive_scale_up_signals = 0

            return new_instances

        elif self.should_scale_down(metrics):
            # Scale down gradually
            new_instances = max(
                self.policy.min_instances,
                self.current_instances - 1
            )

            self.current_instances = new_instances
            self.last_scale_down_time = datetime.now()
            self.consecutive_scale_down_signals = 0

            return new_instances

        return None

    def estimate_cost_savings(
        self,
        without_autoscaling_instances: int,
        cost_per_instance_hour: float,
        hours: int = 24 * 30
    ) -> Dict:
        """Estimate cost savings from autoscaling"""

        # Without autoscaling (always at peak)
        cost_without = without_autoscaling_instances * cost_per_instance_hour * hours

        # With autoscaling (estimate 60% average utilization)
        avg_instances_with_autoscaling = without_autoscaling_instances * 0.6
        cost_with = avg_instances_with_autoscaling * cost_per_instance_hour * hours

        savings = cost_without - cost_with
        savings_percent = (savings / cost_without) * 100

        return {
            "cost_without_autoscaling": cost_without,
            "cost_with_autoscaling": cost_with,
            "monthly_savings": savings,
            "savings_percent": savings_percent
        }

class PredictiveScaler:
    """Predictive autoscaling using historical patterns"""

    def __init__(self, autoscaler: MLAutoscaler):
        self.autoscaler = autoscaler
        self.traffic_patterns: Dict[int, List[float]] = {}  # hour -> [rps]

    def learn_traffic_pattern(self, metrics: ScalingMetrics):
        """Learn traffic patterns over time"""
        hour_of_day = metrics.timestamp.hour

        if hour_of_day not in self.traffic_patterns:
            self.traffic_patterns[hour_of_day] = []

        self.traffic_patterns[hour_of_day].append(metrics.requests_per_second)

        # Keep only last 30 days
        if len(self.traffic_patterns[hour_of_day]) > 30:
            self.traffic_patterns[hour_of_day].pop(0)

    def predict_next_hour_load(self) -> float:
        """Predict load for next hour"""
        next_hour = (datetime.now() + timedelta(hours=1)).hour

        if next_hour not in self.traffic_patterns:
            return 0

        historical_loads = self.traffic_patterns[next_hour]

        if not historical_loads:
            return 0

        # Use 90th percentile for conservative prediction
        return np.percentile(historical_loads, 90)

    def proactive_scale(self) -> Optional[int]:
        """Proactively scale based on prediction"""
        predicted_rps = self.predict_next_hour_load()

        if predicted_rps == 0:
            return None

        # Estimate instances needed
        # Assume 100 RPS per instance
        target_instances = int(np.ceil(predicted_rps / 100))
        target_instances = max(
            self.autoscaler.policy.min_instances,
            min(target_instances, self.autoscaler.policy.max_instances)
        )

        # Only scale up proactively, let reactive scaling handle scale down
        if target_instances > self.autoscaler.current_instances:
            self.autoscaler.current_instances = target_instances
            return target_instances

        return None

# Example usage
if __name__ == "__main__":
    policy = ScalingPolicy(
        min_instances=2,
        max_instances=20,
        target_cpu_utilization=70.0,
        target_latency_p99_ms=200.0
    )

    autoscaler = MLAutoscaler(policy)

    # Simulate high load
    high_load_metrics = ScalingMetrics(
        timestamp=datetime.now(),
        cpu_utilization=85.0,
        memory_utilization=70.0,
        requests_per_second=500,
        avg_latency_ms=150,
        p99_latency_ms=250,
        queue_depth=50
    )

    # Make scaling decision
    new_instances = autoscaler.execute_scaling_decision(high_load_metrics)

    if new_instances:
        print(f"Scaled to {new_instances} instances")

    # Estimate savings
    savings = autoscaler.estimate_cost_savings(
        without_autoscaling_instances=10,
        cost_per_instance_hour=0.50
    )

    print(f"\nAutoscaling Savings:")
    print(f"  Monthly savings: ${savings['monthly_savings']:.2f}")
    print(f"  Savings percent: {savings['savings_percent']:.1f}%")
\end{lstlisting}

\section{Cloud Cost Monitoring and Alerts}

Continuous cost monitoring prevents budget overruns and identifies optimization opportunities.

\subsection{Cost Monitoring System}

\begin{lstlisting}[language=Python, caption=Cloud Cost Monitoring and Alerting]
#!/usr/bin/env python3
"""
Cloud Cost Monitoring
Real-time cost tracking and alerting system
"""

import boto3
from datetime import datetime, timedelta
from typing import Dict, List, Optional
from dataclasses import dataclass
import json

@dataclass
class CostAlert:
    """Cost alert configuration"""
    alert_name: str
    threshold_amount: float
    threshold_percent: Optional[float]
    period_days: int
    notification_emails: List[str]
    severity: str = "warning"  # warning, critical

@dataclass
class CostAnomaly:
    """Detected cost anomaly"""
    service: str
    current_cost: float
    expected_cost: float
    deviation_percent: float
    timestamp: datetime

class CloudCostMonitor:
    """Monitor cloud costs and send alerts"""

    def __init__(self, aws_account_id: str, region: str = "us-east-1"):
        self.account_id = aws_account_id
        self.ce_client = boto3.client('ce', region_name=region)  # Cost Explorer
        self.sns_client = boto3.client('sns', region_name=region)
        self.cloudwatch_client = boto3.client('cloudwatch', region_name=region)
        self.alerts: List[CostAlert] = []

    def get_cost_and_usage(
        self,
        start_date: datetime,
        end_date: datetime,
        granularity: str = "DAILY",
        group_by: Optional[List[Dict]] = None
    ) -> Dict:
        """Get cost and usage data from AWS Cost Explorer"""

        if group_by is None:
            group_by = [{"Type": "DIMENSION", "Key": "SERVICE"}]

        response = self.ce_client.get_cost_and_usage(
            TimePeriod={
                'Start': start_date.strftime('%Y-%m-%d'),
                'End': end_date.strftime('%Y-%m-%d')
            },
            Granularity=granularity,
            Metrics=['UnblendedCost', 'UsageQuantity'],
            GroupBy=group_by
        )

        return response

    def get_current_month_cost(self) -> float:
        """Get month-to-date cost"""

        start_of_month = datetime.now().replace(day=1, hour=0, minute=0, second=0, microsecond=0)
        today = datetime.now()

        response = self.get_cost_and_usage(
            start_date=start_of_month,
            end_date=today,
            granularity="MONTHLY",
            group_by=[]
        )

        if response['ResultsByTime']:
            total_cost = float(
                response['ResultsByTime'][0]['Total']['UnblendedCost']['Amount']
            )
            return total_cost

        return 0.0

    def get_forecast(self, days_ahead: int = 30) -> float:
        """Get cost forecast"""

        start_date = datetime.now()
        end_date = start_date + timedelta(days=days_ahead)

        response = self.ce_client.get_cost_forecast(
            TimePeriod={
                'Start': start_date.strftime('%Y-%m-%d'),
                'End': end_date.strftime('%Y-%m-%d')
            },
            Metric='UNBLENDED_COST',
            Granularity='MONTHLY'
        )

        forecast = float(response['Total']['Amount'])
        return forecast

    def add_cost_alert(self, alert: CostAlert):
        """Add a cost alert"""
        self.alerts.append(alert)

        # Create CloudWatch alarm
        self._create_cloudwatch_alarm(alert)

    def _create_cloudwatch_alarm(self, alert: CostAlert):
        """Create CloudWatch alarm for cost threshold"""

        # Note: AWS Cost Explorer metrics are delayed, so this is simplified
        # In practice, you'd need to set up AWS Budgets or custom metrics

        alarm_name = f"cost-alert-{alert.alert_name}"

        try:
            self.cloudwatch_client.put_metric_alarm(
                AlarmName=alarm_name,
                ComparisonOperator='GreaterThanThreshold',
                EvaluationPeriods=1,
                MetricName='EstimatedCharges',
                Namespace='AWS/Billing',
                Period=86400,  # Daily
                Statistic='Maximum',
                Threshold=alert.threshold_amount,
                ActionsEnabled=True,
                AlarmDescription=f'Cost alert for {alert.alert_name}',
                Dimensions=[
                    {
                        'Name': 'Currency',
                        'Value': 'USD'
                    }
                ]
            )

            print(f"Created CloudWatch alarm: {alarm_name}")

        except Exception as e:
            print(f"Error creating alarm: {e}")

    def check_alerts(self) -> List[str]:
        """Check all alerts and return triggered ones"""

        triggered_alerts = []
        current_cost = self.get_current_month_cost()

        for alert in self.alerts:
            # Calculate threshold for period
            if alert.threshold_percent:
                # Compare to previous period
                prev_period_cost = self._get_previous_period_cost(alert.period_days)
                threshold = prev_period_cost * (1 + alert.threshold_percent / 100)
            else:
                threshold = alert.threshold_amount

            if current_cost > threshold:
                message = (
                    f"Cost Alert: {alert.alert_name}\n"
                    f"Current cost: ${current_cost:.2f}\n"
                    f"Threshold: ${threshold:.2f}\n"
                    f"Severity: {alert.severity}"
                )

                triggered_alerts.append(message)
                self._send_alert_notification(alert, message)

        return triggered_alerts

    def _get_previous_period_cost(self, period_days: int) -> float:
        """Get cost from previous period"""

        end_date = datetime.now() - timedelta(days=period_days)
        start_date = end_date - timedelta(days=period_days)

        response = self.get_cost_and_usage(
            start_date=start_date,
            end_date=end_date,
            granularity="MONTHLY",
            group_by=[]
        )

        if response['ResultsByTime']:
            return float(
                response['ResultsByTime'][0]['Total']['UnblendedCost']['Amount']
            )

        return 0.0

    def _send_alert_notification(self, alert: CostAlert, message: str):
        """Send alert notification via SNS"""

        # In practice, you'd have an SNS topic set up
        print(f"ALERT: {message}")

        # Example SNS publish (requires topic ARN)
        # self.sns_client.publish(
        #     TopicArn='arn:aws:sns:us-east-1:123456789:cost-alerts',
        #     Subject=f'Cost Alert: {alert.alert_name}',
        #     Message=message
        # )

    def detect_cost_anomalies(self) -> List[CostAnomaly]:
        """Detect unusual cost patterns"""

        # Get last 7 days of costs by service
        end_date = datetime.now()
        start_date = end_date - timedelta(days=7)

        response = self.get_cost_and_usage(
            start_date=start_date,
            end_date=end_date,
            granularity="DAILY",
            group_by=[{"Type": "DIMENSION", "Key": "SERVICE"}]
        )

        anomalies = []

        # Analyze each service
        service_costs = {}

        for result in response['ResultsByTime']:
            for group in result['Groups']:
                service = group['Keys'][0]
                cost = float(group['Metrics']['UnblendedCost']['Amount'])

                if service not in service_costs:
                    service_costs[service] = []

                service_costs[service].append(cost)

        # Detect anomalies (simple threshold: >50% above average)
        for service, costs in service_costs.items():
            if len(costs) < 2:
                continue

            avg_cost = sum(costs[:-1]) / len(costs[:-1])
            current_cost = costs[-1]

            if current_cost > avg_cost * 1.5:
                deviation = ((current_cost - avg_cost) / avg_cost) * 100

                anomalies.append(CostAnomaly(
                    service=service,
                    current_cost=current_cost,
                    expected_cost=avg_cost,
                    deviation_percent=deviation,
                    timestamp=datetime.now()
                ))

        return anomalies

    def generate_cost_report(self) -> Dict:
        """Generate comprehensive cost report"""

        current_month_cost = self.get_current_month_cost()
        forecast = self.get_forecast(days_ahead=30)

        # Get cost by service
        start_of_month = datetime.now().replace(day=1, hour=0, minute=0, second=0, microsecond=0)
        today = datetime.now()

        by_service = self.get_cost_and_usage(
            start_date=start_of_month,
            end_date=today,
            granularity="MONTHLY"
        )

        service_costs = {}
        if by_service['ResultsByTime']:
            for group in by_service['ResultsByTime'][0]['Groups']:
                service = group['Keys'][0]
                cost = float(group['Metrics']['UnblendedCost']['Amount'])
                service_costs[service] = cost

        # Detect anomalies
        anomalies = self.detect_cost_anomalies()

        return {
            "current_month_cost": current_month_cost,
            "forecasted_month_end_cost": forecast,
            "cost_by_service": service_costs,
            "anomalies": [
                {
                    "service": a.service,
                    "deviation_percent": a.deviation_percent,
                    "current_cost": a.current_cost
                }
                for a in anomalies
            ],
            "generated_at": datetime.now().isoformat()
        }

    def set_budget(
        self,
        budget_name: str,
        budget_limit: float,
        notification_threshold_percent: float = 80.0
    ):
        """Set up AWS Budget"""

        budgets_client = boto3.client('budgets')

        try:
            budgets_client.create_budget(
                AccountId=self.account_id,
                Budget={
                    'BudgetName': budget_name,
                    'BudgetLimit': {
                        'Amount': str(budget_limit),
                        'Unit': 'USD'
                    },
                    'TimeUnit': 'MONTHLY',
                    'BudgetType': 'COST'
                },
                NotificationsWithSubscribers=[
                    {
                        'Notification': {
                            'NotificationType': 'ACTUAL',
                            'ComparisonOperator': 'GREATER_THAN',
                            'Threshold': notification_threshold_percent,
                            'ThresholdType': 'PERCENTAGE'
                        },
                        'Subscribers': [
                            {
                                'SubscriptionType': 'EMAIL',
                                'Address': 'team@example.com'
                            }
                        ]
                    }
                ]
            )

            print(f"Created budget: {budget_name}")

        except Exception as e:
            print(f"Error creating budget: {e}")

# Example usage
if __name__ == "__main__":
    monitor = CloudCostMonitor(aws_account_id="123456789012")

    # Add cost alerts
    alert = CostAlert(
        alert_name="monthly-budget",
        threshold_amount=10000.0,
        threshold_percent=None,
        period_days=30,
        notification_emails=["team@example.com"],
        severity="warning"
    )

    monitor.add_cost_alert(alert)

    # Generate cost report
    report = monitor.generate_cost_report()

    print(f"Current Month Cost: ${report['current_month_cost']:.2f}")
    print(f"Forecasted Cost: ${report['forecasted_month_end_cost']:.2f}")

    print("\nCost by Service:")
    for service, cost in sorted(
        report['cost_by_service'].items(),
        key=lambda x: x[1],
        reverse=True
    )[:5]:
        print(f"  {service}: ${cost:.2f}")

    if report['anomalies']:
        print("\nCost Anomalies Detected:")
        for anomaly in report['anomalies']:
            print(f"  {anomaly['service']}: +{anomaly['deviation_percent']:.1f}% "
                  f"(${anomaly['current_cost']:.2f})")
\end{lstlisting}

\section{ROI Analysis for ML Investments}

Quantifying the return on ML investments justifies costs and guides resource allocation.

\subsection{ML ROI Calculator}

\begin{lstlisting}[language=Python, caption=ML Investment ROI Calculator]
#!/usr/bin/env python3
"""
ML Investment ROI Calculator
Calculate and analyze return on investment for ML projects
"""

from dataclasses import dataclass
from typing import Dict, List, Optional
from datetime import datetime, timedelta
import numpy as np

@dataclass
class MLProjectCosts:
    """ML project cost breakdown"""
    development_costs: float
    training_infrastructure: float
    serving_infrastructure_monthly: float
    data_costs: float
    personnel_costs_monthly: float
    maintenance_costs_monthly: float
    monitoring_costs_monthly: float

    def total_initial_investment(self) -> float:
        """Calculate total initial investment"""
        return (
            self.development_costs +
            self.training_infrastructure +
            self.data_costs
        )

    def monthly_operating_cost(self) -> float:
        """Calculate monthly operating costs"""
        return (
            self.serving_infrastructure_monthly +
            self.personnel_costs_monthly +
            self.maintenance_costs_monthly +
            self.monitoring_costs_monthly
        )

@dataclass
class MLProjectBenefits:
    """ML project benefits"""
    revenue_increase_monthly: float
    cost_savings_monthly: float
    efficiency_gains_monthly: float
    customer_satisfaction_value_monthly: float

    def total_monthly_benefit(self) -> float:
        """Calculate total monthly benefit"""
        return (
            self.revenue_increase_monthly +
            self.cost_savings_monthly +
            self.efficiency_gains_monthly +
            self.customer_satisfaction_value_monthly
        )

class MLROICalculator:
    """Calculate ROI for ML projects"""

    def __init__(
        self,
        project_costs: MLProjectCosts,
        project_benefits: MLProjectBenefits,
        discount_rate: float = 0.10
    ):
        self.costs = project_costs
        self.benefits = project_benefits
        self.discount_rate = discount_rate

    def calculate_simple_roi(self, months: int = 12) -> float:
        """Calculate simple ROI"""

        total_cost = (
            self.costs.total_initial_investment() +
            self.costs.monthly_operating_cost() * months
        )

        total_benefit = self.benefits.total_monthly_benefit() * months

        roi = ((total_benefit - total_cost) / total_cost) * 100

        return roi

    def calculate_payback_period(self) -> float:
        """Calculate payback period in months"""

        initial_investment = self.costs.total_initial_investment()
        monthly_net_benefit = (
            self.benefits.total_monthly_benefit() -
            self.costs.monthly_operating_cost()
        )

        if monthly_net_benefit <= 0:
            return float('inf')

        payback_months = initial_investment / monthly_net_benefit

        return payback_months

    def calculate_npv(self, months: int = 36) -> float:
        """Calculate Net Present Value"""

        initial_investment = self.costs.total_initial_investment()
        monthly_rate = self.discount_rate / 12

        npv = -initial_investment

        for month in range(1, months + 1):
            monthly_cash_flow = (
                self.benefits.total_monthly_benefit() -
                self.costs.monthly_operating_cost()
            )

            discounted_cash_flow = monthly_cash_flow / ((1 + monthly_rate) ** month)
            npv += discounted_cash_flow

        return npv

    def calculate_irr(self, months: int = 36, precision: float = 0.0001) -> float:
        """Calculate Internal Rate of Return"""

        # Use binary search to find IRR
        low_rate = -0.99
        high_rate = 1.0

        while high_rate - low_rate > precision:
            mid_rate = (low_rate + high_rate) / 2

            # Calculate NPV at mid_rate
            npv = self._npv_at_rate(mid_rate, months)

            if npv > 0:
                low_rate = mid_rate
            else:
                high_rate = mid_rate

        return (low_rate + high_rate) / 2

    def _npv_at_rate(self, annual_rate: float, months: int) -> float:
        """Calculate NPV at a specific discount rate"""

        initial_investment = self.costs.total_initial_investment()
        monthly_rate = annual_rate / 12

        npv = -initial_investment

        for month in range(1, months + 1):
            monthly_cash_flow = (
                self.benefits.total_monthly_benefit() -
                self.costs.monthly_operating_cost()
            )

            discounted_cash_flow = monthly_cash_flow / ((1 + monthly_rate) ** month)
            npv += discounted_cash_flow

        return npv

    def sensitivity_analysis(
        self,
        parameter: str,
        variation_range: tuple = (-0.3, 0.3),
        steps: int = 7
    ) -> Dict:
        """Perform sensitivity analysis on a parameter"""

        variations = np.linspace(variation_range[0], variation_range[1], steps)
        results = {
            "variations": [],
            "roi": [],
            "npv": [],
            "payback_period": []
        }

        original_costs = self.costs
        original_benefits = self.benefits

        for variation in variations:
            # Modify parameter
            if parameter == "revenue_increase":
                self.benefits.revenue_increase_monthly = (
                    original_benefits.revenue_increase_monthly * (1 + variation)
                )
            elif parameter == "serving_cost":
                self.costs.serving_infrastructure_monthly = (
                    original_costs.serving_infrastructure_monthly * (1 + variation)
                )
            # Add more parameters as needed

            # Calculate metrics
            results["variations"].append(variation)
            results["roi"].append(self.calculate_simple_roi())
            results["npv"].append(self.calculate_npv())
            results["payback_period"].append(self.calculate_payback_period())

        # Restore original values
        self.costs = original_costs
        self.benefits = original_benefits

        return results

    def compare_scenarios(
        self,
        scenarios: Dict[str, tuple]
    ) -> Dict:
        """Compare different scenarios"""

        results = {}

        for scenario_name, (costs, benefits) in scenarios.items():
            calculator = MLROICalculator(costs, benefits, self.discount_rate)

            results[scenario_name] = {
                "roi_12_months": calculator.calculate_simple_roi(12),
                "roi_36_months": calculator.calculate_simple_roi(36),
                "payback_period_months": calculator.calculate_payback_period(),
                "npv_36_months": calculator.calculate_npv(36),
                "irr": calculator.calculate_irr(36) * 100  # Convert to percentage
            }

        return results

    def generate_roi_report(self) -> Dict:
        """Generate comprehensive ROI report"""

        roi_12m = self.calculate_simple_roi(12)
        roi_36m = self.calculate_simple_roi(36)
        payback = self.calculate_payback_period()
        npv = self.calculate_npv(36)
        irr = self.calculate_irr(36) * 100

        # Calculate break-even
        monthly_net = (
            self.benefits.total_monthly_benefit() -
            self.costs.monthly_operating_cost()
        )

        # Cost efficiency metrics
        cost_per_prediction = (
            self.costs.monthly_operating_cost() /
            (self.benefits.revenue_increase_monthly / 10)  # Assume $10 per conversion
        )

        return {
            "financial_metrics": {
                "roi_12_months_percent": roi_12m,
                "roi_36_months_percent": roi_36m,
                "payback_period_months": payback,
                "npv_36_months": npv,
                "irr_percent": irr
            },
            "cost_breakdown": {
                "initial_investment": self.costs.total_initial_investment(),
                "monthly_operating_cost": self.costs.monthly_operating_cost(),
                "annual_operating_cost": self.costs.monthly_operating_cost() * 12
            },
            "benefit_breakdown": {
                "monthly_benefit": self.benefits.total_monthly_benefit(),
                "annual_benefit": self.benefits.total_monthly_benefit() * 12
            },
            "efficiency_metrics": {
                "monthly_net_benefit": monthly_net,
                "cost_per_prediction": cost_per_prediction
            },
            "recommendation": self._get_recommendation(roi_36m, payback, npv)
        }

    def _get_recommendation(
        self,
        roi: float,
        payback_months: float,
        npv: float
    ) -> str:
        """Generate investment recommendation"""

        if npv > 0 and roi > 20 and payback_months < 12:
            return "Strongly Recommended: Excellent ROI with quick payback"
        elif npv > 0 and roi > 10 and payback_months < 24:
            return "Recommended: Positive ROI with reasonable payback period"
        elif npv > 0:
            return "Consider: Positive but modest returns"
        else:
            return "Not Recommended: Negative returns expected"

# Example usage
if __name__ == "__main__":
    # Recommendation system project
    costs = MLProjectCosts(
        development_costs=150000,  # 6 months * 2 engineers
        training_infrastructure=5000,
        serving_infrastructure_monthly=3000,
        data_costs=10000,
        personnel_costs_monthly=25000,
        maintenance_costs_monthly=2000,
        monitoring_costs_monthly=500
    )

    # Benefits: 5% increase in conversion rate
    benefits = MLProjectBenefits(
        revenue_increase_monthly=50000,  # Increased sales
        cost_savings_monthly=10000,  # Reduced manual curation
        efficiency_gains_monthly=5000,  # Engineering time saved
        customer_satisfaction_value_monthly=5000  # Reduced churn
    )

    calculator = MLROICalculator(costs, benefits)

    # Generate report
    report = calculator.generate_roi_report()

    print("ML Project ROI Analysis")
    print("=" * 50)
    print(f"\nFinancial Metrics:")
    print(f"  12-month ROI: {report['financial_metrics']['roi_12_months_percent']:.1f}%")
    print(f"  36-month ROI: {report['financial_metrics']['roi_36_months_percent']:.1f}%")
    print(f"  Payback Period: {report['financial_metrics']['payback_period_months']:.1f} months")
    print(f"  NPV (36 months): ${report['financial_metrics']['npv_36_months']:,.2f}")
    print(f"  IRR: {report['financial_metrics']['irr_percent']:.1f}%")

    print(f"\nCost Breakdown:")
    print(f"  Initial Investment: ${report['cost_breakdown']['initial_investment']:,.2f}")
    print(f"  Annual Operating Cost: ${report['cost_breakdown']['annual_operating_cost']:,.2f}")

    print(f"\nBenefit Breakdown:")
    print(f"  Annual Benefit: ${report['benefit_breakdown']['annual_benefit']:,.2f}")

    print(f"\nRecommendation: {report['recommendation']}")
\end{lstlisting}

\section{Best Practices for Cost Optimization}

\begin{bestpractice}
Implement cost allocation tags across all ML resources. Tag by project, team, environment, and cost center to enable detailed cost tracking and chargeback.
\end{bestpractice}

\textbf{Cost Optimization Checklist:}

\begin{enumerate}
    \item \textbf{Training Optimization:}
    \begin{itemize}
        \item Profile training jobs before scaling to expensive hardware
        \item Use mixed precision training (FP16/BF16) for 2-3x speedup
        \item Implement spot instances with checkpointing for 70-90\% savings
        \item Use efficient hyperparameter search (Bayesian, Hyperband)
        \item Cache preprocessed training data
        \item Implement early stopping for unpromising runs
    \end{itemize}

    \item \textbf{Serving Optimization:}
    \begin{itemize}
        \item Right-size serving instances based on actual load
        \item Implement autoscaling with ML-aware policies
        \item Use model compression (quantization, distillation)
        \item Implement response caching for repeated requests
        \item Consider serverless for variable/low traffic
        \item Use batch prediction for non-real-time workloads
    \end{itemize}

    \item \textbf{Storage Optimization:}
    \begin{itemize}
        \item Implement data lifecycle policies (hot/warm/cold tiers)
        \item Compress training data and model artifacts
        \item Delete old checkpoints and experiment artifacts
        \item Use efficient file formats (Parquet, TFRecord)
        \item Archive infrequently accessed data
    \end{itemize}

    \item \textbf{Development Optimization:}
    \begin{itemize}
        \item Automatically shut down unused notebook instances
        \item Use smaller instances for development/testing
        \item Share development infrastructure across teams
        \item Implement resource quotas per user/team
    \end{itemize}

    \item \textbf{Monitoring and Governance:}
    \begin{itemize}
        \item Set up cost alerts and budgets
        \item Review cost reports weekly
        \item Track cost per prediction/request
        \item Calculate and optimize ROI
        \item Implement showback/chargeback for accountability
    \end{itemize}
\end{enumerate}

\section{Chapter Summary}

This chapter covered comprehensive cost optimization strategies for ML systems:

\begin{itemize}
    \item \textbf{Cost Analysis:} Implemented tools to track and analyze ML system costs across training, serving, and storage

    \item \textbf{Training Optimization:} Covered mixed precision, gradient accumulation, spot instances, and efficient hyperparameter search

    \item \textbf{Serving Optimization:} Explored deployment options, autoscaling, model compression, and caching strategies

    \item \textbf{Resource Allocation:} Designed intelligent scaling policies that balance performance and cost

    \item \textbf{Cost Monitoring:} Built monitoring systems with alerts and anomaly detection

    \item \textbf{ROI Analysis:} Developed frameworks to calculate and justify ML investments
\end{itemize}

Cost optimization is not a one-time activity but an ongoing process. Regular analysis, monitoring, and refinement of cost strategies ensure ML systems deliver maximum value relative to their infrastructure investments.

\section{Exercises}

\begin{exercise}
\textbf{Training Cost Analysis}

Profile a deep learning training workload and optimize its cost:

\begin{enumerate}
    \item Profile a model training job on different instance types
    \item Calculate cost per epoch for each instance type
    \item Implement mixed precision training and measure speedup
    \item Calculate cost savings from using spot instances
    \item Determine optimal batch size for cost efficiency
\end{enumerate}

\textbf{Deliverables:}
\begin{itemize}
    \item Performance profiling report with throughput metrics
    \item Cost comparison table across instance types
    \item Recommended configuration with estimated annual savings
\end{itemize}
\end{exercise}

\begin{exercise}
\textbf{Serving Cost Optimization}

Analyze and optimize serving infrastructure costs:

\begin{enumerate}
    \item Collect metrics from a production serving endpoint (RPS, latency, instance utilization)
    \item Compare costs of always-on vs. autoscaling vs. serverless deployment
    \item Implement model compression (quantization or distillation)
    \item Measure latency and throughput impact of compression
    \item Calculate ROI of compression investment
\end{enumerate}

\textbf{Deliverables:}
\begin{itemize}
    \item Current serving cost breakdown
    \item Deployment option comparison with recommendations
    \item Model compression analysis with accuracy/latency/cost tradeoffs
\end{itemize}
\end{exercise}

\begin{exercise}
\textbf{Autoscaling Policy Design}

Design and implement an ML-aware autoscaling policy:

\begin{enumerate}
    \item Analyze historical traffic patterns for a serving endpoint
    \item Design scaling policy based on CPU, latency, and queue depth
    \item Implement predictive scaling using traffic patterns
    \item Simulate autoscaling with historical data
    \item Calculate cost savings compared to fixed capacity
\end{enumerate}

\textbf{Deliverables:}
\begin{itemize}
    \item Traffic pattern analysis and visualization
    \item Autoscaling policy specification
    \item Simulation results showing scaling behavior
    \item Cost savings analysis over 30 days
\end{itemize}
\end{exercise}

\begin{exercise}
\textbf{Cost Monitoring Dashboard}

Build a comprehensive cost monitoring system:

\begin{enumerate}
    \item Set up cost tracking for ML infrastructure (use AWS Cost Explorer API or equivalent)
    \item Create dashboard showing costs by service, project, and environment
    \item Implement cost anomaly detection
    \item Configure budget alerts at 50\%, 80\%, and 100\% thresholds
    \item Generate weekly cost optimization reports
\end{enumerate}

\textbf{Deliverables:}
\begin{itemize}
    \item Cost monitoring dashboard (Grafana/CloudWatch)
    \item Anomaly detection rules and alert configurations
    \item Sample weekly cost optimization report
\end{itemize}
\end{exercise}

\begin{exercise}
\textbf{ML Project ROI Analysis}

Perform ROI analysis for an ML project:

\begin{enumerate}
    \item Identify and quantify all costs (development, infrastructure, personnel)
    \item Quantify benefits (revenue increase, cost savings, efficiency gains)
    \item Calculate ROI, NPV, IRR, and payback period
    \item Perform sensitivity analysis on key assumptions
    \item Create executive summary with investment recommendation
\end{enumerate}

\textbf{Deliverables:}
\begin{itemize}
    \item Detailed cost and benefit breakdown
    \item Financial metrics (ROI, NPV, IRR, payback period)
    \item Sensitivity analysis charts
    \item Executive summary with recommendation
\end{itemize}
\end{exercise}
